\chapter{绪论}

% --- 兼容性宏：若主文档未加载 amsmath/amsfonts，可避免 \eqref/\mathbb 报错 ---
\providecommand{\eqref}[1]{(\ref{#1})}
\providecommand{\mathbb}[1]{\mathbf{#1}}
% ------------------------------------------------------------------------

随着计算机图形学的发展，其研究目标已由早期面向光栅显示和离线渲染的图像合成，逐步扩展为涵盖几何建模、物理仿真、实时交互、智能内容生成和跨域数据融合的综合性技术体系。在这一演进过程中，研究对象不再局限于“如何生成逼真的图像”，而是进一步转向“如何表示、理解、分析和操控三维世界中的几何信息”。相应地，几何也由图形生成中的静态载体，逐渐转变为连接感知、建模、分析、优化、制造与交互的核心媒介。

几何处理（geometry processing）正是在这一背景下发展起来的重要研究方向，其核心任务是对数字几何数据进行表示、分析、重建、优化与编辑。其研究对象已由传统的三角网格扩展到点云、体网格、隐式曲面、距离场以及神经隐式表示等多种形式；其研究方法也由经验性的几何算法发展为融合变分优化、偏微分方程、离散微分几何、统计学习与可微计算的综合框架。由此，几何处理不再只是建模流程中的辅助后处理工具，而是贯穿三维数据获取、结构表达、计算分析与下游应用的基础环节。与此同时，它在诸多应用领域中发挥着关键支撑作用：在影视动画与游戏工业中服务于角色绑定、形变传递和数字资产复用；在工业设计与智能制造中支撑 CAD/CAM/CAE 一体化建模、结构优化与工艺验证；在建筑、医学、机器人、遥感测绘和数字人文等领域，则承担三维重建、形态分析、场景理解、几何修复与可视化表达等任务。由此，数字几何不仅是计算机图形学中的表达对象，也成为连接现实世界数据、计算建模与工程决策的重要中间层。

近年来，随着三维扫描、摄影测量、激光雷达、医学成像、虚拟现实和增强现实等技术的发展，大量复杂对象与场景被持续数字化，并以多种几何形式进入建模、分析、仿真和交互流程，复杂几何处理因此成为计算机图形学中的重要挑战。不断增长的应用需求也推动几何数据的规模和复杂度持续提升（见图~\ref{fig:chap1_complex_mesh_geometries} ）。

%真实场景中的几何对象通常具有高分辨率、细节丰富、拓扑复杂、采样不均、噪声显著和局部结构多样等特点，若直接在这些原始几何对象上施加约束、建立对应、传播属性或控制形变，往往会面临离散结构不规则、数值稳定性不足、全局一致性难以保证以及计算代价高昂等问题。



\begin{figure}[htbp]
  \centering
  \includegraphics[width=0.88\linewidth]{chap1_complex_mesh_geometries}
  \caption[具有高细节密度和复杂局部结构的网格几何示例。]{具有高细节密度和复杂局部结构的网格几何示例~\cite{Maggiordomo2023MicroMesh}。}
  \label{fig:chap1_complex_mesh_geometries}
\end{figure}


\section{几何代理与包围几何}
在复杂几何对象上直接开展分析、重建、优化与编辑，往往受限于原始模型的离散结构。一方面，目标几何常具有拓扑复杂、局部细节丰富、采样分布不均和噪声干扰等特征；另一方面，许多处理任务又要求在全局范围内建立稳定的约束关系，并在边界与内部之间可靠地传播几何信息。因此，若直接在原始对象上完成相关计算和操纵，往往容易出现求解不稳定、结果失真和计算代价过高等问题。

为缓解这一矛盾，几何处理通常不直接在原始模型上工作，而是先为其构造一个更规整、更稳定、也更便于操作的中间表示。根据用途不同，本文将这类中间表示统称为\textbf{几何代理}，并进一步区分为两类：一类主要服务于分析、传递与求解等后续处理，本文称为\textbf{计算代理}；另一类主要服务于交互编辑与形变控制，本文称为\textbf{操纵代理}。

\subsection*{计算代理}
计算代理的核心目标是为复杂几何构造一个更规整、更稳定、也更便于求解的中间域，使投影、属性传递、重网格化乃至 PDE 求解等操作能够在该域上进行。围绕这一目标，已有方法大致可以按照代理作用范围与构造复杂度概括为三类：表面级代理、体计算代理以及近边界薄层代理。

最早一类计算代理主要工作在曲面层面。其基本思想不是直接沿用输入模型原有的离散连接，而是先构造一个更适于计算的替代曲面域，例如公共基域、公共细分或公共三角化，以及内蕴三角化等方法 \cite{praun2001consistent,kraevoy2004cross,JiaoInternationalJO,Sharp2019NavigatingIT,Takayama2022CompatibleIntrinsicTriangulations,Schmidt2023SurfaceMapsViaAdaptiveTriangulations}。在这些方法中，代理可以表现为共享的参数域、两个曲面之间的公共三角化，或与输入曲面保持对应的内蕴离散结构。借助这样的表面级代理，可以更稳定地建立曲面间对应关系与代理计算的任务。然而，这类方法不能处理体问题，难以处理需要在边界与内部之间联合传播信息的情形。

为进一步支持内部属性表达和体域上的数值计算，研究者又提出了体计算代理，例如拓扑一致的体参数化、体公共域以及 polycube 一类规则体域 \cite{Kwok2014DomainConstructionVolumetricCrossParameterization,Tarini2004PolyCubeMaps,Gregson2011AllHexMeshGenerationViaVolumetricPolyCubeDeformation}。与表面级代理相比，这类方法把代理从曲面扩展到实体内部，使内部结构、材料属性和体网格能够在统一的体域中组织和传播。不过，这类方法通常需要全局体分解、模板一致性、锚点约束，或额外构造较为规则的参数域；对于几何复杂、细节丰富的单个模型而言，其构造与优化过程往往较重。在很多只是希望在原始边界附近建立稳定对应并进行局部到全局的信息传播时，这类体计算代理往往显得过于复杂。

介于上述两类方法之间，近边界薄层代理提供了一种兼具两种方法优点的思路。早期 shell map 及其后续曲面壳映射工作，已经提出在输入表面与其 offset surface 之间构造一层壳空间，以承载细节表达和局部映射 \cite{porumbescu2005shell,Jeschke2007InteractiveSmoothAndCurvedShellMapping}。在此基础上，Jiang 等人进一步提出带有双射投影的壳空间框架 \cite{Jiang2020BijectivePI}：该方法在输入曲面外构造由广义棱柱组成的壳，借助棱柱层这一共同参数化域，边界条件、位移、纹理及其他属性可以在输入模型与代理模型之间稳健往返传递，并进一步支持表面 PDE、布尔操作、位移贴图和四面体网格生成等任务。与表面级代理相比，壳空间不再局限于曲面本身，而是在边界附近提供了可用于传播与求解的薄层体域；与体计算代理相比，它又避免了对整个对象进行全局体参数化或规则体域构造的高代价。

这种近边界薄层代理即为\textbf{壳}（shell）。壳结构是具有一定厚度的包围壳体，其几何域由外表面和内表面共同界定，目标对象位于壳体内部，或与其内层边界保持特定的空间关系。壳结构由双层边界围成一个带状区域，因此天然具有从外表面到内表面、再到目标对象附近的空间过渡关系。这种“双层边界—壳体区域”的结构特征，使壳结构适合作为投影、包覆、过渡和双射构造的中间域，为目标对象提供一个可控的外部计算底座。

\begin{figure}[htbp]
  \centering
  \begin{overpic}[width=0.9\linewidth,percent]{figures/chap1_shell_app.pdf}
    \put(28,35){\makebox(0,0)[c]{\small \textbf{(a)} 边界层网格~\cite{aubry2017boundary}}}
    \put(80,35){\makebox(0,0)[c]{\small \textbf{(b)} 复杂纹理~\cite{porumbescu2005shell}}}
    \put(28,3){\makebox(0,0)[c]{\small \textbf{(c)} 衣物模拟~\cite{chen2023multilayerthickshells}}}
    \put(79,3){\makebox(0,0)[c]{\small \textbf{(d)} 图案编辑~\cite{jin2019shell}}}
  \end{overpic}
  \caption{壳结构在几何处理中的典型应用示意。}
  \label{fig:chap1_shell_app}
\end{figure}

\subsection*{操纵代理}
与计算代理不同，操纵代理更强调以较少的控制自由度产生直观、稳定且可预期的几何变化。按照控制元的类型，可以将其概括为点代理、线代理、面代理和体代理：
点代理以离散约束点直接施加位移，适合局部精确编辑；
对于强调线性构件与结构关系的人造模型，曲线代理可借助少量控制曲线传播形变并保持特征连续，对于具有关节层级的模型，骨架代理能够以少量骨骼控制整体姿态与体积变化，典型如骨架驱动形变方法；
表面代理通过控制曲面传递弯折、扭转等形变
体代理通过体域控制整体形状变化。这些代理共同构成了复杂几何由难以直接操纵走向可控编辑的重要桥梁~\cite{Gain2008Survey,Chadwick1989,Singh1998Wires,Borrel1991,Feng1996}。

\begin{figure}[htbp]
  \centering
  \begin{overpic}[width=0.98\linewidth,percent]{figures/chap1_proxy.pdf}
    \put(12,2){\makebox(0,0)[c]{\small \textbf{(a)} 体代理}}
    \put(37,2){\makebox(0,0)[c]{\small \textbf{(b)} 面代理}}
    \put(63,2){\makebox(0,0)[c]{\small \textbf{(c)} 线代理}}
    \put(88,2){\makebox(0,0)[c]{\small \textbf{(d)} 点代理}}
  \end{overpic}
  \caption[控制变形时的四种典型代理：体、面、线、点。]{控制变形时的四种典型代理~\cite{Gain2008Survey}：体、面、线、点。}
  \label{fig:chap1_proxy}
\end{figure}

在包围式操纵代理中，最具代表性的是\textbf{笼}（cage）。笼代理~\cite{Ju2005} 以包围目标对象的粗控制网格作为低维控制层，通过广义重心坐标或其他插值函数，将边界上的位移、旋转和高阶控制量稳定传播到内部高分辨率几何，从而支持交互式建模、角色动画和数字资产复用等任务。进一步地，Harmonic Coordinates 通过保证坐标正性增强了形变的局部性~\cite{Joshi2007}，Green Coordinates 则将笼面法向纳入控制，有助于在大尺度编辑中更好地保持表面细节~\cite{Lipman2008Green}。因此，从本文的分类看，笼是一种典型的操纵代理。

\begin{figure}[htbp]
  \centering
  \begin{overpic}[width=0.9\linewidth,percent]{figures/chap1_cage_app.pdf}
    \put(35,33){\makebox(0,0)[c]{\small \textbf{(a)} 碰撞检测~\cite{calderon2017bounding}}}
    \put(84,33){\makebox(0,0)[c]{\small \textbf{(b)} 模拟代理~\cite{calderon2017bounding}}}
    \put(50,3){\makebox(0,0)[c]{\small \textbf{(c)} 形状编辑~\cite{Thiery2022TriquadGC}}}
  \end{overpic}
  \caption{笼结构在几何处理中的典型应用。}
  \label{fig:chap1_cage_app}
\end{figure}

\subsection*{包围几何}
基于上述两类代理，本文将壳与笼统一定义为\textbf{包围几何}（enclosing geometry）：即从外部包围目标对象，并承载后续映射、属性传播或形变控制的辅助结构。它们的共同点在于都把原始复杂几何转化为一个更规整、更可控的中间域，从而为后续计算提供统一的外部参考边界与内部组织方式。

\begin{figure}[htbp]
  \centering
  \begin{overpic}[width=0.9\linewidth,percent]{figures/chap1_enclosing_geometry.pdf}
    \put(14,2){\makebox(0,0)[c]{\small 目标几何}}
    \put(45,2){\makebox(0,0)[c]{\small 壳结构}}
    \put(79,2){\makebox(0,0)[c]{\small 笼结构}}
  \end{overpic}
  \caption{包围几何中壳结构与笼结构的关系示意。}
  \label{fig:chap1_enclosing_geometry}
\end{figure}

尽管同属包围几何，壳与笼的结构与用途并不相同。壳由内外两层边界共同界定，形成具有厚度的壳体区域，因此更适合作为投影、对应建立和属性传递的计算空间；笼通常仅由单层外边界构成，因此更适合作为由边界控制内部形变的操作界面。换言之，壳更强调“在包围域中计算”，笼更强调“通过包围域操纵”。二者的差异恰好对应了前述计算代理与操纵代理的不同目标，而二者的统一点则在于都以包围域作为复杂几何处理的中间层。

\section{包围几何的构造}

在明确壳与笼分别承担的计算与操纵角色后，进一步的问题便是如何围绕输入对象自动、稳定地构造这类包围几何。总体来看，围绕笼与壳的构造已经积累了大量工作，且相关研究不再只关心“是否能够从外部包住目标”，而是同时关注自动化、紧致包围、流形性、无自交、特征对齐，以及与后续映射、坐标求解和形变控制的兼容性。换言之，当前包围几何构造通常需要同时回答两个问题：一是“能否可靠构造”，二是“构造之后是否好用”。

在笼结构生成方向，相关构造工作已经形成较清晰的技术演进路线，这一点也可由近年的综述工作得到系统概括~\cite{Stroter2024CageSurvey}：从早期的线性多面体笼，逐步发展到结构化的 triquad 笼，再进一步推广到高阶曲边或曲面笼。早期方法多强调以较少面片从外部包围目标对象，而近年的代表性工作则更重视自动化、紧致包围、流形性、无自交、特征对齐以及与后续坐标构造和形变求解的兼容性。Sacht 等人~\cite{Sacht2015} 研究了嵌套笼的稳定构造，\citet{Guo2024RobustCage} 则进一步提出鲁棒粗笼构造框架：先生成满足包围和无自交要求的初始笼，再在受约束重网格化中同时控制复杂度与逼近误差。在线性笼之上，\citet{Thiery2022TriquadGC} 利用 triquad 结构增强了笼边界与语义特征的对齐能力；随后，曲边笼、高阶 Coons 片笼与 B{\'e}zier 曲面片笼又将笼边界从分片线性推广到高阶形式，从而以更少控制量获得更紧、更光顺、也更适于高质量形变的包围边界~\cite{Lin2024PolynomialCC,Qin2024Coons,Xiao2025BezierCage}。这说明笼的构造目标已不再只是“得到一个可包围目标的粗网格”，而是逐步转向“得到一个同时适合控制、插值和编辑的高质量边界”；与此同时，面向三角笼的双调和坐标及其导数也表明，笼边界设计与后续坐标构造正呈现更紧密的一体化趋势~\cite{thiery2024biharmonic}。

在壳结构生成方向，相关工作同样十分丰富，但问题设置更强调厚度控制、偏置稳定性、拓扑合法性以及壳层内部映射的可用性。大体上看，已有方法至少可概括为三类：一类是外包裹或收缩包裹型壳，更强调从复杂输入外侧生成水密且严格包围目标的外壳，以带偏移的 Alpha Wrapping 方法为代表~\cite{Portaneri2022AlphaWrapping}；一类是壳空间或双射壳，更强调在壳层内部建立从一层曲面到另一层曲面的稳定映射，以双射投影壳空间及其高阶推广为代表~\cite{Jiang2020BijectivePI,liu2024curveshell}；另一类则面向显式偏置或等距壳生成，把锐特征保持、复杂度控制和拓扑正确性放在更核心的位置~\cite{Wang2024PCO,Zint2025TopologicalOffsets}。若把最新发展一并考虑，基于平滑环绕数的变分偏置方法进一步说明，壳生成正在从“单纯求一个偏置面”转向几何特征、拓扑合法性与厚度控制的统一优化~\cite{Sun2026VariationalMeshOffsetting}。与仅用于几何逼近的外包层不同，本文更关注第二类壳结构，因为这类壳不仅要从外部包围输入对象，还要在壳层内部提供可计算的过渡空间。\citet{Jiang2020BijectivePI} 系统提出了带有双射投影的壳空间框架，将输入曲面转化为由广义棱柱组成的线性壳，并在壳层内部建立稳定投影，使属性传递、布尔操作和后续数值计算能够在统一空间中进行。在此基础上，\citet{liu2024curveshell} 进一步将线性壳推广到高阶壳，通过更光滑的边界表示、局部操作与约束优化，获得厚度分布更均匀、投影更平滑且更贴合输入边界的壳空间。由此可见，壳的构造同样经历了从“生成可包围的外层结构”到“生成可直接服务于计算的高质量空间代理”的演进。

正因为笼与壳的构造路线已经很多，无论采用哪一条路线，只要包围几何最终需要承担区域组织、映射定义或内部/外部判定的作用，对非水密边界的鲁棒查询都可视为构造包围几何时绕不开的一个共有问题。随着包围边界由分片线性走向曲边和高阶形式，包围域的识别问题不再是一个可以被默认跳过的预处理步骤。对于规则且严格闭合的水密包围几何边界，已有多种基础包含查询方法可以完成内外判定，因此这一问题通常并不构成主要困难。然而，实际应用中的包围几何并非总是水密，尤其在艺术家使用 CAD 软件生成或编辑的曲边形状中，边界常会出现缝隙、破碎或其他脏几何。此时，问题的关键已不再是对理想闭合边界进行一般意义上的内外判断，而是如何在边界不完备的情况下恢复稳定的区域判定能力，并在计算意义上重建可用的包围域。

现有包含查询方法虽然种类较多，但在破碎曲边边界场景下往往难以同时兼顾鲁棒性、精度与效率。一类方法依赖射线相交、局部线性化或其他离散判定策略，其结果容易受到边界缝隙、退化构型与数值误差的影响；另一类方法通过离散近似或自适应细分处理曲边边界，虽然能够在一定程度上提高几何精度，但计算代价往往会随曲线形状和查询点位置显著波动，难以提供稳定高效的查询性能。因此，仍有必要构造一种面向破碎包围几何的高效鲁棒包含查询方法，以为后续计算提供可靠的区域识别基础。只有在借助鲁棒包含查询稳定识别由破碎包围几何所描述的区域之后，包围几何作为中间计算域的前提才真正成立；在此基础上，后续需要解决的主体问题，才是在包围几何内构造稳定、有效且质量足够高的映射。基于这一点，下一节进一步讨论包围几何内映射问题及其质量需求。
\section{包围几何内的映射}

在包围域能够被直接识别，或经由鲁棒包含查询修复后可被稳定识别的前提下，包围几何中的映射主要涉及两类问题：

\begin{figure}[htbp]
  \centering
  \begin{overpic}[width=0.9\linewidth,percent]{figures/chap1_geometry_mapping.pdf}
    \put(14,2){\makebox(0,0)[c]{\small 目标几何}}
    \put(47,2){\makebox(0,0)[c]{\small 壳空间中的映射}}
    \put(82,2){\makebox(0,0)[c]{\small 笼空间中的映射}}
  \end{overpic}
  \caption{包围几何内壳空间映射与笼空间映射。}
  \label{fig:chap1_geometry_mapping}
\end{figure}

\begin{enumerate}
  \item \textbf{壳空间中的映射。} 以壳结构的边界曲面或壳体区域为定义域，直接构造向量值映射 $\mathbf{f}$。这类问题更关注壳体区域中的几何过渡、边界对应以及映射的双射性与光滑性。
  \item \textbf{笼空间中的映射。} 以笼结构围成的控制域为定义域，构造坐标函数 $\{w_i\}$ 并由其诱导映射。这类问题更关注边界控制量向内部的稳定传播，以及插值性、局部性与形变质量。
\end{enumerate}

为便于讨论，本文将这两类问题统一表述为如下函数构造问题：在域 $\Omega \subset \mathbb{R}^d$ 上构造向量值映射
\begin{equation}
  \mathbf{f}: \Omega \rightarrow \mathbb{R}^d,
\end{equation}
或构造一组标量坐标函数 $\{w_i:\Omega\rightarrow\mathbb{R}\}$，并通过这些坐标函数对边界控制量进行插值，从而得到相应映射。前者主要对应壳空间中的映射，后者主要对应笼空间中的映射。
在离散层面，$\Omega$ 可以由壳面的三角网格、壳体区域的体网格，或笼内部的采样点集表示，$\mathbf{f}$ 或 $\{w_i\}$ 则在相应的顶点或单元上离散。
\subsection*{映射质量与应用需求}

映射质量不仅是理论上的性质，也直接决定属性传递与几何编辑能否稳定进行。

\begin{itemize}
  \item 几何属性传递。
  若能在源域与目标域之间构造稳定映射 $\mathbf{f}$，则颜色、纹理坐标、法向、位移、材料参数等属性即可通过推前或拉回（push-forward/pull-back）传递到目标模型。
  在这一过程中，映射的双射性与畸变水平直接影响结果是否出现重叠、撕裂或模糊。
  \item 几何编辑。
  在笼编辑中，坐标函数 $\{w_i\}$ 的光滑性决定形变的连续性与视觉自然度；
  共形性、低扭曲与局部性则关系到细节保持和局部控制能力；
  若映射在编辑过程中发生折叠，还可能引发自交与不可逆形变，进而影响后续动画与仿真。
\end{itemize}

因此，本文关注的并非仅仅是“能否构造一个映射”，而是“能否构造一个\emph{可直接服务于下游任务}的高质量映射”。包含查询的意义也主要在于为这一映射问题补足可用包围域这一基础前提。
\section{包围几何内的映射构造目标}

本文所讨论的“高质量映射”并非单一指标，而是一组与下游任务直接相关的性质。对包围几何内的映射而言，问题的关键不只是“是否存在一个映射”，而是能否在满足应用需求的前提下，同时兼顾几何质量、边界控制与数值可计算性。
%综合前述讨论，映射构造的核心目标可概括为：
%
%\begin{enumerate}
%  \item \textbf{双射性与可逆性：} 对包围几何空间中的映射，应在离散意义下尽量保证局部 Jacobian 不翻转（$\det J_{\mathbf{f}} > 0$），并在可能情况下获得全局单射，从而避免折叠与自交。
%  \item \textbf{光滑性：} 相比仅有 $C^0$ 连续性的映射，具有 $C^1/C^2$ 连续性的映射能够显著减少折痕与尖峰伪影，对形变传播以及法向、曲率等高阶属性的稳定传递尤为重要。
%  \item \textbf{共形性与低扭曲：} 根据任务需求，映射应尽量保持局部角度结构，并控制剪切、拉伸与扭曲，以避免局部形状在传递或形变过程中发生明显失真。
%  \item \textbf{边界控制与约束一致性：} 映射或坐标应准确响应边界给定的控制条件。对笼坐标而言，这通常表现为严格的边界插值（控制点处 Kronecker Delta 性质） 以及稳定的局部控制能力。
%  \item \textbf{效率与可复现：} 方法应能在常见网格规模下高效求解，最好具有可预分解或可重用的线性系统结构，并给出清晰的离散化与实现细节。
%  \item \textbf{可验证性：} 方法应能对映射质量给出数值证据或证书，例如翻转检测、畸变上界估计以及条件数/稳定性分析，以降低隐藏失效的风险。
%\end{enumerate}
各种质量目标在壳空间映射与笼空间映射中具有不同侧重，因此有必要分别分析。

\subsection{壳空间中的映射：双射性与光滑性}

壳空间中的映射可以看作一层层面之间的映射关系，例如外层与内层之间的对应，或沿壳体厚度方向的过渡映射。这类映射首先承担投影、包覆和边界传递的作用，因此其核心性质集中在双射性与光滑性上。

\begin{itemize}
  \item \textbf{双射性。} 若层与层之间的映射失去单射性，壳空间中便会出现折叠、自交或覆盖混乱，进而破坏投影关系和后续属性传递的正确性。因此，壳空间映射首先需要尽量保证局部不翻转，并在可能情况下保持全局可逆。
  \item \textbf{光滑性。} 壳空间映射不仅要建立对应关系，还要为后续的内部传播提供稳定边界条件。若映射不够光滑，局部几何误差会在传递过程中被放大，影响属性场、法向以及形变结果的连续性。
\end{itemize}

基于这一点，本文在壳结构上将问题聚焦为光滑双射映射的构造，并围绕生成与映射相容的壳空间展开研究。

然而，这一目标在离散壳结构上并不容易实现。一方面，许多现有方法虽然能够在几何上相近的曲面之间建立对应关系，但由于几何的复杂性，往往难以稳定保证映射的双射性。另一方面，即便能够获得双射映射，仍然难以进一步构造一个光滑映射，以往的壳空间结构难以与这样的映射相容。为了服务属性传递、重网格化和后续内部传播，壳空间不仅需要在层间维持稳定的一一对应，还需要使投影在壳内连续变化，并避免局部空间过窄、分布不均或边界处不光滑所带来的数值问题。因此，如何同时兼顾双射性、光滑性以及壳空间的可构造性，构成了壳结构上映射设计的核心挑战。

\subsection{笼空间中的映射：共形性、扭曲控制与边界控制}

笼空间中的映射本质上是在笼所围成的体域内传播边界控制量，因此可以看作一种体域上的映射或坐标传播机制。与壳空间中的层间映射
不同，这类映射需要将边界上的位移、旋转等控制信息稳定地传递到整个内部区域。其可逆性通常主要依赖用户对笼边界的控制；当笼边
界在编辑过程中保持不自交时，所诱导的内部映射通常也更容易保持可逆。在此基础上，笼空间映射还需要兼顾光滑性，以获得平滑自然
的变形效果，并进一步关注共形性、边界控制以及扭曲控制等性质。

\begin{itemize}
	\item \textbf{共形性。} 当映射尽量保持局部角度结构时，内部形变通常更加自然，局部细节也不易因角度畸变而失真。
	\item \textbf{边界控制。} 笼方法的核心优势在于由边界控制内部，因此映射必须准确响应边界条件，保持严格边界插值以
	及可预期的局部控制能力，否则控制点的几何意义将被削弱。
	\item \textbf{扭曲控制。} 严格的边界控制与共形性之间是相互排斥的：除非该映射本身就是仿射变换，否则不能同时满足
	共形与插值要求。因此，在强调边界控制的情形下，对共形性的要求往往需要转化为对整体扭曲水平的控制。若笼空间映射在体域
	内部产生过大的剪切、拉伸或扭曲，就容易导致不自然的体积变形和局部伪影，因此，控制扭曲同样是保证形变质量的关键要求。
\end{itemize}

许多经典坐标可以理解为某类 PDE 的边值问题解，例如调和坐标与双调和坐标。除此之外，也存在基于复分析与共形结构构造的坐标体系。围绕上述两类思想，本文在笼结构上沿两条路线展开。

\begin{enumerate}
  \item \textbf{高阶柯西坐标。}
  柯西坐标是一类拥有共形性质的坐标，为了使其拥有光滑性，本文想要构造一类高阶柯西坐标，即推广柯西坐标适用的笼的边界的阶数，采用由更光滑参数曲线构成的笼边界来承载形变控制。
  
    目前已有许多算法研究如何计算高阶柯西坐标，然而，之前的算法要么是对光滑性的小范围改进，即在给定条件下是光滑的，要么采取了极为复杂的复合映射与数值方法。这不仅使坐标与导数难以获得统一的闭式表达，也对输入笼施加了额外几何限制，从而削弱了笼在表达能力、编辑直观性与计算效率上的优势。
  \item \textbf{高阶双调和坐标。}
  双调合坐标是拥有边界控制和扭曲控制能力的坐标，严格的边界控制也表明其变形结果的光滑性需要笼边界的光滑性，为此，本文想要构造一类高阶双调合坐标。
  
  然而，将双调和坐标从线性笼推广到高阶情形并不直接。其困难不仅在于高阶曲边边界下相关积分计算更为复杂，也在于边值问题的建立、离散方程系统的构造以及坐标求解过程都会随边界表示的提升而显著复杂化。因此，如何为高阶双调和坐标建立统一的边值问题与离散求解框架，并在此基础上实现半解析计算，是这一方向中的关键问题。
\end{enumerate}

由此可见，壳空间映射与笼空间映射虽然都服务于包围几何上的高质量映射，但它们关注的核心性质并不相同：前者重点在于双射性与光滑性，后者则在双射性和光滑性的基础上进一步强调共形性、扭曲控制与边界控制。下一节将在此基础上概括本文围绕这些目标展开的四项具体工作。
\section{本文内容与结构安排}

在上一节明确映射构造目标后，本节进一步概括本文的主要内容与整体结构。整体框架如图~\ref{fig:chap1_structure} 所示。本文的核心主题是包围几何内的高质量映射。围绕这一主题，本文首先讨论当包围几何因缝隙、破碎或脏几何而无法被直接稳定识别时，如何通过鲁棒内外判断恢复可用的包围域；在此前提得到修复后，再分别研究壳结构与笼结构内的高质量映射。具体包含四项核心工作：

\begin{figure}[htbp]
  \centering
  \includegraphics[width=0.92\linewidth]{chap1_structure.png}
  \caption{文章结构。}
  \label{fig:chap1_structure}
\end{figure}

\begin{enumerate}
  \item \textbf{面向破碎包围几何的高效鲁棒的内外判断。}
  本文研究由有理参数曲线构成的二维包围边界上的鲁棒内外判断问题。方法上，本文以广义环绕数作为统一判据，将有理参数曲线上的内外判断转化为多项式 B\'{e}zier 曲线上的有理多项式积分，并借助留数定理给出闭式计算表达；同时进一步处理查询点位于边界上的特殊情形，并推导广义环绕数导数。由此得到的方法可在复杂曲边边界下实现高效、稳定的内外判定，并在计算意义上修复存在缝隙或破碎的包围边界，为后续高质量映射与坐标构造恢复稳定的区域判定基础。
  \item \textbf{壳空间上的光滑双射投影。}
  本文研究壳空间内的光滑双射投影问题，提出一种高阶壳表示及其实用构造算法。该方法以初始线性壳为起点，通过边塌缩、几何优化、边翻转与棱柱缩放等局部操作，在满足包围约束、角度约束和厚度约束的前提下，逐步得到由高阶三角棱柱组成的壳结构，并在壳内构造连续向量场以定义投影算子。实验表明，该方法能够稳定生成高阶壳，并获得更光滑、更均匀的壳空间投影，适合支撑重网格化后的属性回传、布尔运算与位移贴图等应用。
  \item \textbf{笼空间上的高阶柯西坐标。}
  本文研究高阶笼上的 Cauchy 坐标，给出坐标、导数与相关有理积分的统一闭式解析表达。针对传统线性输入笼对复杂边界贴合不足、边界附近光滑性欠佳的问题，本文将 Cauchy 坐标推广到高阶输入笼，并基于留数定理构建有理积分的闭式求解流程。实验结果表明，高阶笼能够在保持保角变形优势的同时提供更贴合边界、更平滑且更符合编辑意图的结果。
  \item \textbf{笼空间上的高阶双调和坐标。}
  本文研究高阶笼上的双调和坐标，提出面向曲边边界的高阶 BEM 求解框架。该方法通过高阶边界单元、高阶形函数与统一代数系统，将双调和边值问题推广到高阶笼，并结合不同边界导数设计与加权策略，在严格边界插值与内部低扭曲之间提供连续可调的平衡机制。实验结果表明，本文方法能够在复杂曲边笼编辑中提供更稳定的内部形变传播，并在变形自然性与边界可控性之间取得更灵活的折中。
\end{enumerate}

在上述四项工作中，第一项工作的作用是恢复可用的包围域，从而为后三项高质量映射与坐标构造奠定前提；后三项工作则分别从壳空间投影、笼空间共形坐标与低扭曲坐标三个角度展开本文的主体研究。在此基础上，本文进一步讨论它们在几何属性传递、几何编辑与包含查询中的应用。

除本章外，全文其余章节安排如下：

\begin{itemize}
  \item 第2章介绍相关理论基础与研究现状，主要包括几何映射的离散表示、畸变度量、广义重心坐标以及基于 PDE 的坐标构造等内容。
  \item 第3章围绕破碎包围几何上的鲁棒包含查询展开，以有理参数曲线构成的二维包围边界为对象，讨论广义环绕数的闭式计算及其效率表现，说明如何为存在缝隙或脏几何的包围边界恢复稳定的区域判定基础，并为后续以包围几何为计算域的映射与坐标构造补足前提条件。
  \item 第4章研究壳结构上的光滑双射投影模型、算法与性质，并通过属性传递实验验证方法的有效性。
  \item 第5章介绍高阶柯西坐标的构造方法，讨论其数学性质、离散实现以及在笼编辑中的表现。
  \item 第6章从第5章的共形坐标出发，转向研究边界可控且内部低扭曲的高阶双调和坐标，重点讨论其半解析半数值构造、变形控制以及与经典坐标的对比实验。
  \item 第7章总结全文工作，并讨论后续可进一步开展的研究方向。
\end{itemize}
